% !TeX root = thesis.tex
\noindent Predicting gene expression responses after drug treatment is difficult because a drug can affect more than its main target, and the result also depends on the starting state of the cell. This thesis studies whether known biological information can improve this prediction task.

\noindent The proposed model is Uncertainty Perturbation (UPert). It combines a matched control expression profile, drug target information, and a gene interaction graph from OmniPath. The prediction target is the response vector, defined as the difference between treatment and matched control expression. The model predicts both a mean response and an uncertainty value for each gene, which makes it possible to study how reliable the predictions are.

\noindent Experiments are performed on LINCS L1000 landmark genes under three test settings: warm, cold drug target, and cold cell. UPert is compared with two baseline models, DeepCOP and GSNN, under the same split and pairing framework. In these experiments, UPert gives the most consistent results across the main metrics. It performs best on PCC averaged across samples and top gene precision in all three splits, and it also gives the lowest MSE in warm and cold drug target. In cold cell, DeepCOP gives a slightly lower MSE. Cold cell is the hardest setting on PCC and top gene precision.

\noindent The uncertainty analysis shows that higher predicted uncertainty is usually linked to larger prediction error, especially in harder cases. Attention case studies also show that highly weighted edges can be useful clues when predictions are accurate, but these weights should not be treated as direct biological proof.
